\chapter{Preparar el entorno y los datos}
\label{cap:preparar}

\begin{enpocas}
Comprobar que PyTorch ve la tarjeta gráfica, conocer el dataset de fuego y humo
que genera el simulador y crear la carpeta donde vivirá la red.
\end{enpocas}

\paso{1} \textbf{Comprobar el entorno de Python.}
La simulación ya tiene un entorno virtual para YOLO en \cod{~/.venvs/yolo}, con
PyTorch, OpenCV y Ultralytics. FuegoNet solo necesita PyTorch, OpenCV, NumPy y
PyYAML; Ultralytics se usa únicamente para comparar.

\begin{lstlisting}[style=bash]
~/.venvs/yolo/bin/python -c "import torch, cv2, numpy; \
  print(torch.__version__, torch.cuda.is_available(), torch.cuda.get_device_name(0)); \
  print(cv2.__version__, numpy.__version__)"
\end{lstlisting}

\begin{resultado}
\ttfamily\small 2.14.1+cu130 True NVIDIA GeForce RTX 5060 Ti\\
4.11.0 1.26.4
\end{resultado}

\begin{aviso}
NumPy debe ser de la serie 1.x en este entorno: el puente de ROS 2 Jazzy
(\cod{cv_bridge}, \cod{rclpy}) está compilado contra NumPy~1. Si instalas algo con
\cod{pip} y te sube NumPy a 2.x, regrésalo de inmediato:
\begin{lstlisting}[style=bash]
~/.venvs/yolo/bin/pip install "numpy==1.26.4" "ml-dtypes<0.6"
\end{lstlisting}
Nos pasó al instalar \cod{onnx}: pip subió NumPy a 2.5 sin avisar
(capítulo~\ref{cap:integrar}).
\end{aviso}

\paso{2} \textbf{Conocer el dataset.}
Las imágenes salen del propio simulador. \cod{scripts/capturar_dataset.py}
vuela el dron sobre los mundos \cod{bosque} y \cod{bosque_grande} con focos de
fuego en lugares al azar y escribe, para cada foto, un archivo de texto con las
cajas. Son dos datasets que se usan juntos:

\begin{table}[H]
\centering\small
\caption{Dataset de fuego y humo (imágenes de $640\times480$).}
\label{tab:dataset}
\begin{tabular}{llrrrr}
\toprule
\textbf{Dataset} & \textbf{Parte} & \textbf{Imágenes} & \textbf{Cajas fuego} & \textbf{Cajas humo} & \textbf{Sin objetos}\\
\midrule
\cod{incendio_bosque} & train & 2\,700 & 1\,927 & 2\,203 & 309\\
                      & val   &   300 &   213 &   239 &  40\\
\cod{incendio_bosque_grande} & train & 2\,700 & 891 & 1\,898 & 755\\
                      & val   &   300 &    98 &   213 &  81\\
\midrule
\textbf{Total} & train & 5\,400 & 2\,818 & 4\,101 & 1\,064\\
               & val   &   600 &   311 &   452 &   121\\
\bottomrule
\end{tabular}
\end{table}

Las imágenes «sin objetos» son importantes: le enseñan a la red que un claro
iluminado o una roca oscura \emph{no} son fuego (una roca dio falsos positivos al
principio; por eso se agregó el mundo grande).

El formato es el de YOLO. Cada imagen \cod{images/train/00001.jpg} tiene su
\cod{labels/train/00001.txt} con una línea por objeto:
\begin{lstlisting}
0 0.472461 0.706724 0.070710 0.106154
1 0.480463 0.596074 0.145035 0.212378
\end{lstlisting}
El primer número es la clase (0 = fuego, 1 = humo) y los otros cuatro son el
centro $(c_x, c_y)$, el ancho $w$ y el alto $h$ de la caja, \emph{divididos} entre
el ancho y el alto de la imagen (todo queda entre 0 y 1). Un archivo vacío
significa «no hay nada que detectar».

Los dos datasets se juntan con un \cod{data.yaml} que lista ambas carpetas:
\begin{lstlisting}
train:
  - /home/jairo/datasets/incendio_bosque/images/train
  - /home/jairo/datasets/incendio_bosque_grande/images/train
val:
  - /home/jairo/datasets/incendio_bosque/images/val
  - /home/jairo/datasets/incendio_bosque_grande/images/val
names:
  0: fuego
  1: humo
\end{lstlisting}
Está en \cod{~/datasets/combinado_incendio_bosque_incendio_bosque_grande/}
(lo creó \cod{scripts/entrenar_yolo.py}).

\begin{tecnico}
De las 6\,919 cajas de entrenamiento, 195 tienen un lado menor de 16 píxeles.
Son focos vistos desde lejos. Sin un cuidado especial (STAL,
capítulo~\ref{cap:asignacion}) esas cajas pueden quedar sin ninguna celda
que aprenda de ellas.
\end{tecnico}

\paso{3} \textbf{Crear la carpeta de la red.}
FuegoNet vive dentro del paquete de ROS 2 como un subpaquete de Python, para que
el nodo detector la pueda importar después:

\begin{lstlisting}[style=bash]
cd ~/ros2_ws/src/cuadrirotor_sim
mkdir -p cuadrirotor_sim/red_fuego
touch cuadrirotor_sim/red_fuego/__init__.py
\end{lstlisting}

Al terminar el manual la carpeta tendrá estos archivos:
\begin{lstlisting}
cuadrirotor_sim/red_fuego/
  __init__.py        bloques.py      modelo.py      asignador.py
  perdida_iou.py     perdida.py      datos.py       optimizador.py
  metricas.py        entrenar.py     evaluar.py     inferir.py
  exportar.py
\end{lstlisting}

Como son módulos de un paquete, se ejecutan con \cod{python -m} desde la
carpeta del paquete (así funcionan los \cod{import} relativos como
\cod{from .bloques import Conv}).
